\begin{frame}{Rappel: notion de schéma}
\begin{block}{Fondamentale en XML~:}
    \begin{itemize}
        \item Mettre en place un vocabulaire contrôlé;
        \item Contraindre la structure d'un document (et l'adapter à nos besoins).
    \end{itemize}
    Vérifiable avec un validateur XML. Ex: \url{https://www.xmlvalidation.com/}.
\end{block}
\begin{block}{}
    Pour les contraintes en TEI, par exemple:
    \begin{itemize}
        \item \og utilisez \texttt{<placeName>} pour les noms de lieux, \texttt{<persName>} pour les noms de personnes et \texttt{<name>} pour tout autre nom\fg
    \end{itemize}
    \textrightarrow{} Ne peut pas être vérifié manuellement (on regarde les guidelines).
\end{block}
\end{frame}

\begin{frame}{Pourquoi contraindre la TEI?}
    Très pratique dans les projets d'édition numérique permet par exemple de~:
    \begin{itemize}
        \item Contraindre les valeurs de l'attribut \texttt{@type};
        \item Documenter les choix d'édition.
    \end{itemize}
\end{frame}

\begin{frame}{Les schémas XML}
    \begin{itemize}
        \item DTD;
        \item XML Schema: crée par le W3C, écrit en XML, permet de définir la structure et le type de contenu d'un document XML;
        \item Relax NG: écrit en XML, supporte les espaces de noms XML, peut créer des contraintes précises.
    \end{itemize}
\end{frame}

\begin{frame}[fragile]{Les schémas XML: DTD}
    \begin{lstlisting}[numbers=none]
    <!ELEMENT texte (chapitre+)>
    <!ELEMENT chapitre (titre?, paragraphe+)>
    <!ELEMENT titre (#PCDATA)>
    <!ATTTLIST chapitre CDATA #REQUIRED>
    \end{lstlisting}
\end{frame}

\begin{frame}[fragile]{Les schémas XML: XML Schema}
    \begin{lstlisting}[numbers=none]
    <xs:element name="texte" maxOccurs="unbounded">
         <xs:complexType>
             <xs:sequence>
             <xs:element name="titre"
             type="xs:string"/>
             <xs:element name="chap" type="xs:decimal"
             minOccurs="0"/>
             </xs:sequence>
         </xs:complexType>
    </xs:element>
    \end{lstlisting}
\end{frame}

\begin{frame}[fragile]{Les schémas XML: Relax NG}
    \begin{lstlisting}[numbers=none]
    <element name="edition">
         <zeroOrMore>
             <element name="texte">
                 <element name="titre">
                    <text/>
                 </element>
                 <element name="paragraphe">
                    <text/>
                 </element>
             </element>
         </zeroOrMore>
    </element>
    \end{lstlisting}
\end{frame}

\begin{frame}{Combinaison de la DTD et du RNG: l'ODD}
\begin{block}{\textit{One Document Does It All}}
\vspace{0.5em}
    Comporte à la fois \textbf{un schéma de validation} et \textbf{une documentation} de ce schéma.\\
    \textrightarrow{} un seul document XML avec les informations pour le traitement informatique et pour l'humain.
\end{block}
\begin{block}{Deux types de personnalisation basiques}
\begin{itemize}
    \item \textbf{\textit{tei\_all} =} autorise tous les éléments définis par la TEI.
    \item \textbf{\textit{tei\_lite} =} propose environ une cinquantaine d'élément TEI (au lieu de 450).
\end{itemize}

\end{block}
    
\end{frame}

\begin{frame}[fragile]{Structure ODD: déclarations}
    \begin{lstlisting}
        <?xml version="1.0" encoding="UTF-8"?>
        <?xml-model href="url vers odd rng" 
        type="application/xml" schematypens="http://relaxng.org/ns/structure/1.0"?>
        <?xml-model 
        href="url vers odd rng"
        schematypens="http://purl.oclc.org/dsdl/schematron"?>
        
    \end{lstlisting}
\end{frame}
\begin{frame}[fragile]{Structure ODD: en-tête}
    \begin{lstlisting}
        <TEI xmlns="http://www.tei-c.org/ns/1.0" xmlns:xi="http://www.w3.org/2001/XInclude"
          xmlns:sch="http://purl.oclc.org/dsdl/schematron">
          <teiHeader>
            <fileDesc>
              <titleStmt></titleStmt>
              <publicationStmt></publicationStmt>
              <sourceDesc></sourceDesc>
            </fileDesc>
            <revisionDesc></revisionDesc>
          </teiHeader>
          <text>
            <body></body>
          </text>
        </TEI>
    \end{lstlisting}
\end{frame}

\begin{frame}[fragile]{Structure ODD: corps du texte}
    \begin{itemize}
        \item \texttt{<schemaSpec>:} défini les spécifications du schéma;
        \item \texttt{<moduleRef>:} défini \href{https://www.tei-c.org/Vault/P5/1.5.0/doc/tei-p5-doc/fr/html/REF-CLASSES-MODEL.html}{un module} qui doit être présent dans l'encodage;
        \item \texttt{<elementSpec>:} permet de modifier, ajouter, contraindre un élément;
        \item \texttt{<attList>:} permet de contraindre le contenu d'un attribut.
    \end{itemize}
\end{frame}

\begin{frame}{Structure ODD: corps du texte}
    \begin{figure}
        \centering
        \includegraphics[width=\linewidth]{02 - Structurer un document XML-TEI/exemple_moduleref.png}
    \end{figure}
\end{frame}

\begin{frame}{Structure ODD: corps du texte}
    \begin{figure}
        \centering
        \includegraphics[width=0.9\linewidth]{02 - Structurer un document XML-TEI/exemple_attlist.png}
    \end{figure}
\end{frame}

\begin{frame}{Créer son ODD}
   \begin{itemize}
       \item Outil Roma: \url{https://roma.tei-c.org/members}
   \end{itemize}
\end{frame}

\begin{frame}{Exercice: générer une ODD pour le \textit{Bourgeois Gentilhomme}}
Générer une ODD en utilisant l'outil Roma.
    
\end{frame}